\section{Discussion}\label{sec:discussion}

% This section currently contains mechanistic interpretation of the per-area
% weighted EA results from the drug target validation. Additional discussion
% paragraphs (e.g., model limitations, broader implications) can be added here.

\subsection{Therapeutic Area--Specific Performance of Indirect Support}

% DISCUSSION PARAGRAPH 1: Oncology -- literature saturation defeats OTP
% This paragraph explains WHY OTP anti-predicts oncology drug success.
The per-area analysis of weighted EA reveals that the relative performance of PIGEAN and literature-based resources varies systematically with the structure of each therapeutic area's druggable landscape. In oncology, where PIGEAN indirect support achieved the highest weighted EA of any area--method combination ($\text{EA}|\text{RS}{=}1$ = {{WEA_PIGEAN_INDIRECT_ONCO:,}}; $N$ = {{WEA_PIGEAN_INDIRECT_ONCO_N:,}}) and OTP combined was negative ({{WEA_OTP_COMBINED_ONCO:,}}; $N$ = {{WEA_OTP_COMBINED_ONCO_N:,}}; Figure~\ref{fig:weighted-ea}), the failure of literature-based scoring reflects a signal saturation problem. OTP assigns near-identical combined scores (0.95--0.999) to thousands of oncology gene--disease pairs because the cancer literature is vast: nearly every gene has been studied in the context of cancer, compressing OTP's score distribution and eliminating discriminative power. Consistent with this, OTP's top-half and bottom-half oncology targets have identical launched rates (7.1\% versus 7.1\%), meaning OTP's ranking provides zero predictive signal for which targets succeed clinically. PIGEAN's pathway-based indirect support, by contrast, scores genes based on their functional biology rather than publication frequency: PIGEAN's top-half oncology targets have a launched rate of 7.8\%, compared with 4.1\% for the bottom half, indicating genuine ranking discrimination. Among PIGEAN's top-ranked launched oncology targets are TP53 (head and neck carcinoma), SMO (basal cell carcinoma), JAK2 (polycythemia vera), MAP2K1 and MAP2K2 (thyroid carcinoma), BCL2 (chronic lymphocytic leukemia), and VEGFA (renal cell carcinoma), all targets identified through pathway membership rather than literature volume.

% DISCUSSION PARAGRAPH 2: Neurology -- cross-phenotype pathway integration
In neurology, where PIGEAN combined support achieved six-fold higher weighted EA than OTP ({{WEA_PIGEAN_COMBINED_NEURO:,}} versus {{WEA_OTP_COMBINED_NEURO}}; Figure~\ref{fig:weighted-ea}), the top gene sets driving indirect support reveal a coherent mechanistic basis for PIGEAN's advantage. The highest-contributing gene sets include amyloid-$\beta$ metabolic process (summed $\beta$ = 7.68), amyloid precursor protein metabolism (7.10), and $\alpha$-synuclein inclusion body (3.49), representing the core pathobiology of Alzheimer's and Parkinson's disease, alongside neuromuscular synapse morphology (7.64), ciliopathies (10.71), cilium assembly (6.71), mitochondrial DNA metabolic process (5.03), and Huntington's disease (2.07). This cross-phenotype integration is the distinguishing feature of indirect support: it borrows strength from related biological annotations, including neurodegenerative, neuromuscular, ciliary, and mitochondrial pathways, to prioritize targets for specific neurological indications. Literature co-occurrence, which scores genes independently for each disease term, cannot perform this kind of integrative reasoning across related phenotypes.

% DISCUSSION PARAGRAPH 3: Cardiovascular -- indirect support outperforms combined
The cardiovascular result, where PIGEAN indirect support (weighted EA = {{WEA_PIGEAN_INDIRECT_CARDIO}}; $N$ = {{WEA_PIGEAN_INDIRECT_CARDIO_N:,}}) exceeded combined support ({{WEA_PIGEAN_COMBINED_CARDIO}}; $N$ = {{WEA_PIGEAN_COMBINED_CARDIO_N:,}}; Figure~\ref{fig:weighted-ea}), provides direct evidence that pathway-based reasoning can outperform direct genetic association for drug target identification. The targets uniquely prioritized by indirect support, but not by combined support, are predominantly ion channel and calcium-handling genes with strong pathway biology but weak or absent GWAS signals: RYR2 for cardiac arrhythmias (indirect = 5.16, direct = 0.0), LCAT for myocardial infarction (indirect = 4.92, direct = 0.01), CACNA1C and CACNA1D for atrial fibrillation (indirect = 4.54 and 4.45; direct = 0.31 and 0.0), and KCNH2 for arrhythmias (indirect = 4.16, direct = 0.0), a launched antiarrhythmic target identified entirely by indirect evidence. The top gene sets for cardiovascular disease, including contractile fiber (summed $\beta$ = 7.83), arrhythmogenic right ventricular cardiomyopathy (7.75), cardiac conduction (5.56), cardiac fibrosis (5.68), cardiac hypertrophy (5.13), cholesterol metabolism (5.47), and plasma lipoprotein assembly and clearance (5.18), directly reflect the mechanistic pharmacology of cardiovascular drugs. Adding direct GWAS support dilutes this pathway signal by introducing genes with statistically significant but biologically ambiguous associations that do not map to established druggable biology.

% DISCUSSION PARAGRAPH 4: Immune -- OTP's literature depth captures a well-studied space
The immunology result presents an informative contrast: OTP indirect support achieved weighted EA = {{WEA_OTP_INDIRECT_IMMUNE}} ($N$ = {{WEA_OTP_INDIRECT_IMMUNE_N:,}}), while PIGEAN indirect support was near zero ({{WEA_PIGEAN_INDIRECT_IMMUNE}}; $N$ = {{WEA_PIGEAN_INDIRECT_IMMUNE_N:,}}; Figure~\ref{fig:weighted-ea}). OTP's top immune targets are the canonical cytokine and receptor drug targets: TNF (juvenile arthritis, launched), IL6 (rheumatoid arthritis, launched), IL17A (juvenile arthritis, launched), SERPING1 (hereditary angioedema, launched), and IGHE (asthma, Phase III), all thoroughly characterized in the immunology literature. PIGEAN's top immune targets, by contrast, are T-cell surface markers and signaling molecules (CD28, CD6, CD19, IL7R, CD200, LCK, CD3E) that are biologically central to immune function but whose clinical translation has been limited by safety and selectivity challenges. Only 6 of the top 50 genes ranked by PIGEAN and OTP overlapped. This pattern is consistent with the hypothesis that immunology drug development has been dominated by targeting well-characterized soluble mediators for which decades of literature faithfully record clinical translatability, whereas PIGEAN's pathway-based inference identifies central immune biology that is harder to exploit pharmacologically.

% DISCUSSION PARAGRAPH 5: Metabolic -- complementary ranking of lipid vs. glucose targets
In metabolic disease, OTP combined and PIGEAN combined achieved nearly identical weighted EA above RS = 1 ({{WEA_OTP_COMBINED_METABOLIC:,}} versus {{WEA_PIGEAN_COMBINED_METABOLIC:,}}; Figure~\ref{fig:weighted-ea}), but the $\text{EA}|\text{RS}{=}\text{Self}$ metric revealed that PIGEAN provides substantially better ranking within its target set ({{WEA_SELF_PIGEAN_COMBINED_METABOLIC:,}} versus {{WEA_SELF_OTP_COMBINED_METABOLIC}}). The ranking differences reflect the underlying evidence sources. PIGEAN assigns its highest scores to lipid metabolism targets: INS for type 2 diabetes (rank 1, combined = 12.3, launched), LDLR for hyperlipidemia (rank 2, combined = 11.8), PCSK9 for atherosclerosis (rank 69, combined = 10.0, launched), and NPC1L1 for hyperlipidemia (rank 77, combined = 9.87, launched), all targets supported by richly annotated cholesterol and lipoprotein pathway gene sets. OTP preferentially ranks glucose and incretin pathway targets: SLC5A2 for type 2 diabetes (OTP rank 14 versus PIGEAN rank 2,770, launched), GLP1R (OTP rank 49 versus PIGEAN rank 595, launched), and DPP4 (OTP rank 145 versus PIGEAN rank 1,810, launched). The discrepancy is consistent with the relative annotation depth of these pathways in MSigDB and the Mammalian Phenotype Ontology: lipid metabolism is extensively annotated across dozens of gene sets (cholesterol metabolism, lipoprotein pathways, sterol biosynthesis), whereas glucose transport and incretin signaling have sparser pathway representation. This finding suggests that the performance of indirect support is bounded by the completeness of the underlying gene set libraries, and that expanding annotation coverage of under-represented biological pathways could further improve PIGEAN's drug target prioritization.
